\documentclass[12pt]{article}
\usepackage[utf8]{inputenc}
\usepackage[T1]{fontenc}
\usepackage[margin=1in]{geometry}
\usepackage{setspace}
\usepackage{xcolor}
\usepackage[colorlinks=true, urlcolor=blue!60!black, linkcolor=blue!60!black]{hyperref}
\setlength{\parindent}{2em}
\setlength{\parskip}{1.25em}
\renewcommand{\baselinestretch}{1.0}
\begin{document}
\begin{center}
{\LARGE \textit{Trump v. Anderson}, 601 U.S. 100 (2024)\par}
\vspace{0.8em}
{\large Case Brief --- Four-Page Edition\par}
\end{center}
\vspace{1em}

\section{Memory Jogger}

Colorado's Supreme Court removed Donald Trump from the 2024 presidential primary ballot under Section~3 of the Fourteenth Amendment; the Court unanimously reversed, holding that States may disqualify state officers and candidates but have no power to enforce Section~3 against federal offices, which a five-Justice majority assigned to Congress (\href{https://www.supremecourt.gov/opinions/23pdf/601us1r06_m6hn.pdf}{601 U.S. 100}).

\section{Detailed Case Facts}

In September 2023, six months before Colorado's March 5, 2024 presidential primary, four Republican and two unaffiliated Colorado voters, led by Norma Anderson, petitioned a Colorado state court under the Colorado Election Code against former President Donald J. Trump and Secretary of State Jena Griswold. They alleged that after losing the 2020 election, Trump intentionally organized and incited the crowd that breached the U.S. Capitol on January 6, 2021, disqualifying himself from the Presidency under Section~3.

\section{Procedural History}

After a five-day trial, Denver District Judge Sarah B. Wallace found that Trump had ``engaged in insurrection'' but denied the petition because the President is not an ``officer of the United States'' covered by Section~3. On December 19, 2023, the Colorado Supreme Court, 4--3, held that Section~3 covers the Presidency, is self-executing and justiciable under the Election Code, and barred Trump from the primary ballot; Chief Justice Boatright and Justices Samour and Berkenkotter dissented. The Supreme Court granted expedited certiorari on January 5, heard argument February 8, and decided the case on March 4, 2024, one day before Super Tuesday.

\section{Judicial Votes}

Reversed, \textbf{9--0} in the judgment, dividing \textbf{5--4} on whether congressional legislation is the exclusive means of enforcing Section~3.
\begin{itemize}
\item \textbf{Per curiam:} Joined by Chief Justice Roberts and Justices Thomas, Alito, Gorsuch, and Kavanaugh.
\item \textbf{Concurrences:} Justice Barrett (in part and in the judgment); Justices Sotomayor, Kagan, and Jackson (jointly, in the judgment).
\item \textbf{Dissents:} None.
\end{itemize}

\section{Holding}

All nine Justices held that States lack power to enforce Section~3 against federal officeholders and candidates, especially presidential candidates, though they may disqualify persons holding or seeking state office. A five-Justice majority further held that responsibility for enforcing Section~3 against federal officeholders and candidates rests with Congress under Section~5, through tailored legislation. The Court did not decide whether Trump engaged in insurrection or whether the President is an ``officer of the United States.''

\section{Analysis of Opinions}

\textbf{Per curiam.} \textit{Federalism.} The Fourteenth Amendment ``expand[ed] federal power at the expense of state autonomy.'' States may set qualifications for their own officers and enforced Section~3 against them after ratification, but power over federal officers must be delegated, not reserved (\textit{U.S. Term Limits}). \textit{Text and amnesty.} Nothing delegates enforcement over federal candidates to States, and pre-election exclusion would burden Congress's power to remove a disability by two-thirds of each House at any time, contrary to \textit{McCulloch}. \textit{History and Section~5.} Citing \textit{Griffin's Case} and Senator Trumbull, the majority stressed that a ``preventive and severe penalty'' requires formal procedures, which Congress supplied in the Enforcement Act of 1870 and the Confiscation Act; States, unbound by \textit{City of Boerne}'s limits, could not enjoy freer rein than Congress. \textit{Patchwork.} Presidential elections implicate a ``uniquely important national interest''; divergent state standards would yield conflicting outcomes and post-election chaos.

\textbf{Barrett (concurring in part).} She declined to decide whether federal legislation is the exclusive vehicle, urged that the Court ``turn the national temperature down,'' and stressed the Justices' unanimity on the outcome.

\textbf{Sotomayor, Kagan, and Jackson (concurring in the judgment).} Invoking Chief Justice Roberts's restraint in \textit{Dobbs}, they faulted the majority for deciding how federal enforcement must work. Section~3 is self-executing like other Reconstruction provisions; requiring a simple-majority statute undercuts the two-thirds amnesty rule; \textit{Griffin's Case} is a nonprecedential circuit decision; and the majority sought to ``insulate this Court and petitioner from future controversy.''

\section{Examples: Future Applications}

\begin{enumerate}
\item \textbf{Same-Side (No State Power): Senate Candidate.} A State excludes a former Senator from its Senate primary for aiding a rebellion. A Senate seat is federal, so the State lacks power, whatever the Senate's own Article~I, \S~5 authority.
\item \textbf{Same-Side (No State Power): Sitting Representative.} A state court removes an incumbent Representative who helped plan an armed attack on a federal building. The insurrection finding does not matter.
\item \textbf{Opposite-Side (State Power): Governor's Race.} On the same record, the State bars a former sheriff from the gubernatorial ballot, within the carve-out for state offices.
\item \textbf{Opposite-Side (Likely Valid): Underage Candidate.} A 29-year-old is kept off a presidential ballot. Age is objective and unwaivable by Congress, as in \textit{Hassan v.\ Colorado}.
\item \textbf{Fence-Sitter: Two-Term Limit.} A State excludes a twice-elected former President under the Twenty-Second Amendment. The limit is objective, but \textit{Anderson}'s concern about state-by-state presidential decisions may still apply.
\end{enumerate}

\section{Critique}

\textbf{Textual inversion and selective originalism.} Requiring affirmative legislation lets a bare majority, or one chamber's inaction, nullify Section~3, though its text requires a two-thirds supermajority of both Houses to lift a disability. Critics also note that the Court relied on consequentialist ``patchwork'' fears and a single lower-court precedent rather than original public meaning.

\textbf{Unresolved merits.} The Court left open whether the President is an ``officer of the United States,'' what conduct is ``insurrection,'' and whether Congress could reject electoral votes on Section~3 grounds in joint session without prior legislation.

\section{Key Quotations}

\begin{quote}
``We conclude that States may disqualify persons holding or attempting to hold state office. But States have no power under the Constitution to enforce Section 3 with respect to federal offices, especially the Presidency.''
\end{quote}
\noindent --- Per Curiam, 601 U.S. at 110.

\begin{quote}
``Because the Constitution makes Congress, rather than the States, responsible for enforcing Section 3 against federal officeholders and candidates, we reverse.''
\end{quote}
\noindent --- Per Curiam, 601 U.S. at 106.

\begin{quote}
``The Court today needed to resolve only a single question: whether an individual State may keep a Presidential candidate found to have engaged in insurrection off its ballot. The majority resolves much more than the case before us.''
\end{quote}
\noindent --- Justices Sotomayor, Kagan, and Jackson, concurring in the judgment, 601 U.S. at 123.

\end{document}
